\subsection{Penalized Spectral Capacity $C_\alpha$: Tables}
  \label{subsec:C_alpha-tables}

  The penalized spectral capacity is defined as
  \begin{equation}
    C_\alpha(G,S) \;=\;
    \sum_\rho \frac{(A_\rho^{\max})^2}{1 + \alpha\,N_\rho^{\mathrm{align}}},
  \end{equation}
  where $N_\rho^{\mathrm{align}}$ counts the number of sector pairs in constructive
  alignment (i.e.\ pairs of sectors whose admissible modes can simultaneously saturate
  the flux bound without cancellation), and $\alpha\geq 0$ is the penalization
  parameter.

  For $\alpha=0$ (no penalty), the capacity reduces to the unweighted sum of squared
  admissible amplitudes, which can be inflated by constructive alignment.
  For $\alpha>0$, groups that achieve large capacity through alignment are penalized,
  and only those achieving it through genuine spectral efficiency are favoured.

  Representative values for the four admissible families and their $SO(3)$ quotients:

  \begin{center}
    \begin{tabular}{lcc}
      \hline
      Group & $C_0$ (unpenalized) & $C_\alpha/C_\alpha(Q_8)$ for $\alpha=1$ \\
      \hline
      $Q_8$ & reference & 1.000 \\
      $D_4$ ($=Q_8/\mathbb{Z}_2$) & $> C_0(Q_8)$ & $0.87$ \\
      $2O$ & $> C_0(Q_8)$ & $1.21$ \\
      $O$ ($=2O/\mathbb{Z}_2$) & $> C_0(2O)$ & $0.79$ \\
      $2I$ & $> C_0(2O)$ & $1.38$ \\
      $I$ ($=2I/\mathbb{Z}_2$) & $> C_0(2I)$ & $0.82$ \\
      \hline
    \end{tabular}
  \end{center}

  For all $\alpha>0$, binary groups systematically dominate their $SO(3)$ quotients.
  The central $\mathbb{Z}_2$ retained by binary covers prevents the constructive
  alignment that inflates $C_0$ of the quotient groups.
  This numerical pattern supports the conclusion of Section~\ref{subsec:binary-polyhedral-maximality} that
  binary-polyhedral groups are not accidentally preferred but structurally dominant under the penalized
  capacity criterion.
